% !TeX root = Lec01_Diff_Equ.tex

\ifdefined\AdvMacroMain
  \chapter{Difference Equation}
  \let\ThisNoteEnd\relax
\else
  \documentclass[12pt]{article}
  \newcommand{\ProjectRoot}{..}
  \usepackage[a4paper,margin=0.85in]{geometry}
\usepackage{newpxtext}
\usepackage{amsmath,amssymb,amsthm,mathtools}
\usepackage{newpxmath}
\usepackage{xcolor}
\usepackage{enumitem}
\usepackage{graphicx}
\usepackage{array,tabularx}
\usepackage{float}
\usepackage{hyperref}

\definecolor{noteRed}{RGB}{190,45,90}

\newcommand{\red}[1]{\textcolor{noteRed}{#1}}
\newcommand{\blue}[1]{\textcolor{blue}{#1}}
\newcommand{\purple}[1]{\textcolor{purple}{#1}}

\newcommand{\E}{\mathbb{E}}
\newcommand{\dd}{\mathrm{d}}
\newcommand{\MPN}{\mathrm{MPN}}
\newcommand{\MRS}{\mathrm{MRS}}
\newcommand{\MRT}{\mathrm{MRT}}
\newcommand{\Var}{\operatorname{Var}}
\newcommand{\boxit}[1]{\fbox{\ensuremath{#1}}}
\newcommand{\circled}[1]{\textcircled{\scriptsize #1}}

\newenvironment{tightalign}
  {\[\begin{aligned}}
  {\end{aligned}\]}

\graphicspath{
  {\ProjectRoot/_assets/figures/}
  {\ProjectRoot/_assets/figures/lecture_04/}
}

  \title{Difference Equation}
  \author{Chengyuan He}
  \date{March 2026}
  \begin{document}
  \maketitle
  \def\ThisNoteEnd{\end{document}}
\fi

\section{Difference operators}

\[
\text{1-st order difference:}\qquad \Delta y_t = y_t - y_{t-1}
\]

\[
\text{2-nd order difference:}\qquad
\Delta^2 y_t
= \Delta\left(\Delta y_t\right)
= \Delta\left(y_t-y_{t-1}\right)
= \Delta y_t - \Delta y_{t-1}
= \left(y_t-y_{t-1}\right)-\left(y_{t-1}-y_{t-2}\right)
\]
\[
= y_t - 2y_{t-1} + y_{t-2}.
\]

\section{Period}

Period: an interval of time $\Rightarrow y_t$ is the $y$ for $t^{\text{th}}$ interval of time, not for a point of time.

\textcolor{red}{* $t$ is a period concept, not a point-in-time concept.}

\[
\Delta y_t = y_t - y_{t-1}
\]
is the change from $\left(t-1\right)^{\text{th}}$ time interval to $t^{\text{th}}$ interval, not a derivative at a point of time $t$.

\textcolor{red}{* $\Delta y_t$ is a change, not a derivative or growth rate at time $t$.}

\[
\text{Derivative} \;=\; \dot y_t
\;=\;
\frac{\Delta y_t}{y_{t-1}}
\;=\;
\frac{y_t-y_{t-1}}{y_{t-1}}
\;=\;
\frac{y_t}{y_{t-1}}-1
\;\approx\;
\log\!\left(\frac{y_t}{y_{t-1}}\right)
\]

\[
\textcolor{red}{
\text{when } \frac{y_t-y_{t-1}}{y_{t-1}} \text{ is small, }
\log\!\left(1+\frac{y_t}{y_{t-1}}-1\right)
\approx
\frac{y_t}{y_{t-1}}-1
=
\frac{y_t-y_{t-1}}{y_{t-1}}.
}
\]

\section{Solve a first-order difference equation}

Solve a first-order difference equation $\Longleftrightarrow$ rewrite $y_t$ as a function of $t$.

Example:
\[
y_t + a y_{t-1} = C.
\]

\subsection{Brute-force iterative method}

\[
y_t = -a y_{t-1} + C
\]
\[
y_{t-1} = -a y_{t-2} + C
\]

Hence
\[
y_t = a^2 y_{t-2} - aC + C,
\qquad
y_{t-2} = -a y_{t-3} + C.
\]

Continuing,
\[
y_t = -a^3 y_{t-3} + a^2 C - aC + C - \cdots
\]

and in general,
\[
y_t = \left(-a\right)^t y_0 + C \sum_{j=0}^{t-1} \left(-a\right)^j.
\]

If $a\neq -1$,
\[
\sum_{j=0}^{t-1} \left(-a\right)^j
=
\frac{1-\left(-a\right)^t}{1+a}.
\]

So
\[
y_t
=
\left(-a\right)^t\left(y_0-\frac{C}{1+a}\right)+\frac{C}{1+a},
\qquad a\neq -1.
\]

Equivalently,
\[
y_t = A\left(-a\right)^t + \frac{C}{1+a},
\qquad a\neq -1.
\]

If $a=-1$,
\[
\sum_{j=0}^{t-1}\left(-a\right)^j = \sum_{j=0}^{t-1}1=t
\qquad\Rightarrow\qquad
y_t = y_0 + Ct.
\]

Thus
\[
y_t=
\begin{cases}
A\left(-a\right)^t + \dfrac{C}{1+a}, & a\neq -1,\\[0.8em]
Ct+y_0, & a=-1,
\end{cases}
\qquad
\text{where } A = y_0-\dfrac{C}{1+a}.
\]

\subsection{General method}

Since the system is \textit{linear}, we can separate:
\begin{enumerate}[label=\arabic*)]
    \item the part driven by constant $C$,
    \item the part comes from the dynamics.
\end{enumerate}

\[
\text{Why } y_t=y_t^p+y_t^c \text{ is valid?}
\]

Define linear operator
\[
L\left(\cdot\right)= \left(\cdot\right)_t + a\left(\cdot\right)_{t-1}.
\]

Then
\[
y_t + ay_{t-1} = C
\qquad\Rightarrow\qquad
L\left(y\right)=C.
\]

If $y^p$ is any particular solution such that
\[
L\left(y^p\right)=C,
\]
and if $y^c$ is any solution of the homogeneous equation
\[
L\left(y^c\right)=0
\qquad
\text{(complementary solution)},
\]
then by linearity,
\[
L\left(y^p+y^c\right)=L\left(y^p\right)+L\left(y^c\right)=C+0=C.
\]

Therefore,
\[
y^p+y^c
\]
solves the equation
\[
y_t+ay_{t-1}=C.
\]

\[
\boxed{\text{All solutions}=\text{one particular solution}+\text{all homogeneous solutions}.}
\]

\textcolor{red}{Now we solve the difference equation with general method.}

\subsubsection{Complementary solution}
Since exponential sequences are eigenfunctions of the lag operator, try
\[
y_t = Ab^t.
\]

Then for the homogeneous equation,
\[
Ab^t + aAb^{t-1}=0
\quad\Rightarrow\quad
b+a=0
\quad\Rightarrow\quad
b=-a.
\]

Hence
\[
y_t^c = A\left(-a\right)^t.
\]

\subsubsection{Particular solution}
Since the forcing term is a constant $C$, we guess
\[
y_t^p = k.
\]

Then
\[
k+ak=C
\quad\Rightarrow\quad
k=\frac{C}{1+a},
\qquad a\neq -1.
\]

If $a=-1$,
\[
y_t-y_{t-1}=C
\quad\Rightarrow\quad
y_t^p=Ct.
\]

Therefore,
\[
y_t = y_t^c+y_t^p
=
\begin{cases}
A\left(-a\right)^t+\dfrac{C}{1+a}, & a\neq -1,\\[0.8em]
A+Ct, & a=-1.
\end{cases}
\]

Using the initial condition:

If $a\neq -1$,
\[
y_0=A\left(-a\right)^0+\frac{C}{1+a}
=
A+\frac{C}{1+a}
\quad\Rightarrow\quad
A=y_0-\frac{C}{1+a}.
\]

If $a=-1$,
\[
y_0=A
\quad\Rightarrow\quad
A=y_0.
\]

\section{Dynamic stability of equilibrium}

Consider
\[
y_t + ay_{t-1}=C.
\]

The steady state is defined by
\[
y_t=y_{t-1}=\bar y,
\]
so
\[
\bar y + a\bar y = C
\quad\Rightarrow\quad
\bar y=\frac{C}{1+a}
\qquad \left(a\neq -1\right).
\]

Subtract steady state from the original equation:
\[
y_t-\bar y + a\left(y_{t-1}-\bar y\right)
=
C-\bar y-a\bar y.
\]

Define deviation from steady state as
\[
\tilde y_t = y_t-\bar y.
\]

Then
\[
\tilde y_t + a\tilde y_{t-1}=0
\quad\Rightarrow\quad
\tilde y_t = \left(-a\right)\tilde y_{t-1}.
\]

Let
\[
b=-a.
\]
Then
\[
\tilde y_t = b\tilde y_{t-1}
\quad\Rightarrow\quad
\tilde y_t = b^t \tilde y_0.
\]

\textcolor{red}{* $b=-a$ controls convergence.}

\[
\begin{cases}
\left|b\right|<1 \Rightarrow b^t\to 0,\; \tilde y_t\to 0,\; y_t\to \bar y \quad \text{(convergence)},\\
\left|b\right|>1 \Rightarrow b^t\to \infty,\; \tilde y_t\to \infty \quad \text{(divergence)}.
\end{cases}
\]

\section{\texorpdfstring{Sign of $b$ and oscillation / monotonicity}{Sign of b and oscillation / monotonicity}}

Sign of $b$ also controls \textcolor{red}{oscillation} / \textcolor{red}{monotonicity}.

\[
\begin{cases}
b>0:
& b^t \text{ keeps the same sign for all } t
\Rightarrow
\tilde y_t \text{ keeps the same sign for all } t,\\
& \text{path is non-oscillatory (monotone, towards / away from } \bar y).\\[0.5em]
b<0:
& b^t \text{ alternates sign }
\Rightarrow
\tilde y_t \text{ flips sign each period} \\
& \Rightarrow \text{oscillate around } \bar y.
\end{cases}
\]

Further:
\[
\begin{cases}
b>1, & y_t \text{ monotonically diverges from } \bar y,\\
b=1, & y_t \text{ keeps the same distance away from } \bar y,\\
0<b<1, & y_t \text{ monotonically converge to } \bar y,\\
b=0, & y_t \text{ stays at steady state } \bar y \text{ always},\\
-1<b<0, & y_t \text{ oscillate around and converge to } \bar y,\\
b=-1, & y_t \text{ oscillate around } \bar y \text{ and keeps same absolute distance},\\
b<-1, & y_t \text{ oscillate around and diverge from } \bar y.
\end{cases}
\]

\begin{center}
    \IfFileExists{\ProjectRoot/_assets/figures/relation.png}{%
        \includegraphics[width=0.5\linewidth]{relation.png}%
    }{%
        \fbox{\texttt{missing figure: relation.png}}%
    }
\end{center}

\section{\texorpdfstring{The role of $A$}{The role of A}}

Since the general solution can be written as
\[
y_t=\bar y + Ab^t
\quad\Rightarrow\quad
A=y_0-\bar y
\]
(initial deviation).

Scale effect:
\[
\left|A\right| \uparrow \Rightarrow \text{start farther from equilibrium} \Rightarrow \text{bigger movements along path.}
\]

Mirror effect:
\[
A>0:\ \text{start above } \bar y;
\qquad
A<0:\ \text{start below } \bar y.
\]

\section{Example: Inventory model}

\textcolor{red}{
Example: Inventory model.
}

\[
\textcolor{red}{
Q_{d,t}=\alpha-\beta p_t \qquad \text{demand}
}
\]

\[
\textcolor{red}{
Q_{s,t}=-\gamma+\delta p_t \qquad \text{supply}
}
\]

\[
\textcolor{red}{
p_{t+1}=p_t-\sigma\,\left(Q_{s,t}-Q_{d,t}\right)
}
\]

\[
\textcolor{red}{
\alpha,\beta,\gamma,\delta,\sigma>0
\qquad
\text{solve for } p_t\text{'s path}
}
\]

Then
\[
Q_{s,t}-Q_{d,t}
=
-\left(\alpha+\gamma\right)+\left(\beta+\delta\right)p_t.
\]

Plug in:
\[
p_{t+1}
=
\left[1-\sigma\left(\beta+\delta\right)\right]p_t
+
\sigma\left(\alpha+\gamma\right).
\]

At steady state,
\[
p_{t+1}=p_t=\bar p
\quad\Rightarrow\quad
\bar p=\frac{\alpha+\gamma}{\beta+\delta}.
\]

Thus, the deviation dynamics is
\[
p_t-\bar p
=
\left[1-\sigma\left(\beta+\delta\right)\right]\left(p_{t-1}-\bar p\right).
\]

Stability condition tells us:
\[
\left|1-\sigma\left(\beta+\delta\right)\right|<1
\quad\Rightarrow\quad
0<\sigma\left(\beta+\delta\right)<2,
\]
system converges.

Oscillation condition tells us:
\[
1-\sigma\left(\beta+\delta\right)<0
\quad\Rightarrow\quad
\sigma\left(\beta+\delta\right)>1
\]
creates oscillation.

\textcolor{red}{
* Economic intuition: if price adjustment is too aggressive $\sigma \uparrow\uparrow$, creates oscillation and divergence.
}

If supply or demand elasticity in terms of price too high
\[
\beta \uparrow\uparrow / \delta \uparrow\uparrow
\]
\[
\Rightarrow \text{creates oscillation or even divergence.}
\]

\section{Higher-order difference equations}

\[
\text{2-nd order linear difference equation with constant coefficients \& constant term:}
\]
\[
y_{t+2}+a_1 y_{t+1}+a_2 y_t = C.
\]

\subsection{\texorpdfstring{Particular solution $y^p$: the intertemporal equilibrium level}{Particular solution yp: the intertemporal equilibrium level}}

Try a constant:
\[
y^p=k
\quad\Rightarrow\quad
k\left(1+a_1+a_2\right)=C
\quad\Rightarrow\quad
k=\frac{C}{1+a_1+a_2},
\qquad 1+a_1+a_2\neq 0.
\]

Try a first-degree polynomial:
\[
y^p=kt.
\]

Then
\[
k\left(t+2\right)+a_1k\left(t+1\right)+a_2kt = C
\]
\[
\Rightarrow\quad
k\left[\left(1+a_1+a_2\right)t+\left(2+a_1\right)\right]=C.
\]

Impose
\[
1+a_1+a_2=0
\]
then
\[
k=\frac{C}{2+a_1},
\qquad 2+a_1\neq 0.
\]

Try a second-degree polynomial:
\[
y^p=kt^2.
\]

Then
\[
k\left(t+2\right)^2 + a_1 k\left(t+1\right)^2 + a_2 kt^2 = C.
\]

The notes impose
\[
\begin{cases}
1+a_1+a_2=0,\\
2+a_1=0,
\end{cases}
\qquad\Rightarrow\qquad
\begin{cases}
a_1=-2,\\
a_2=1.
\end{cases}
\]

Substituting:
\[
k\left(t^2+4t+4\right)-2k\left(t^2+2t+1\right)+kt^2=C
\]
\[
\Rightarrow\quad 2k=C
\quad\Rightarrow\quad
k=\frac{C}{2}.
\]

Therefore
\[
y^p=
\begin{cases}
\dfrac{C}{1+a_1+a_2}, & 1+a_1+a_2\neq 0,\\[1em]
\dfrac{C}{2+a_1}\, t, & 1+a_1+a_2=0,\ a_1\neq -2,\\[1em]
\dfrac{C}{2}\, t^2, & a_1=-2,\ a_2=1.
\end{cases}
\]

\subsection{\texorpdfstring{Complementary solution $y^c$: deviation from the equilibrium in each period}{Complementary solution yc: deviation from the equilibrium in each period}}

Look at the homogeneous equation:
\[
y_{t+2}+a_1y_{t+1}+a_2y_t=0.
\]

Try
\[
y_t=Ab^t.
\]

Then
\[
b^2+a_1b+a_2=0.
\]

\subsubsection{Two distinct real roots}
If
\[
a_1^2-4a_2>0,
\]
then
\[
y^c=A_1b_1^t+A_2b_2^t
\]
where
\[
b_1=\frac{-a_1+\sqrt{a_1^2-4a_2}}{2},
\qquad
b_2=\frac{-a_1-\sqrt{a_1^2-4a_2}}{2}.
\]

\subsubsection{Repeated real root}
If
\[
a_1^2-4a_2=0,
\]
then
\[
b=b_1=b_2=-\frac{a_1}{2},
\]
and
\[
y^c=A_3 b^t + A_4 t b^t.
\]

\subsubsection{Complex roots}
If
\[
a_1^2-4a_2<0,
\]
then
\[
b_{1,2}=h\pm vi,
\qquad
h=-\frac{a_1}{2},
\qquad
v=\frac{\sqrt{4a_2-a_1^2}}{2}.
\]

So
\[
y^c=A_5\left(h+vi\right)^t+A_6\left(h-vi\right)^t.
\]

Using De Moivre's Theorem:
\[
\left(h\pm vi\right)^t = R^t\left(\cos\left(\theta t\right)\pm i\sin\left(\theta t\right)\right)
\]
where
\[
R=\sqrt{h^2+v^2}=\sqrt{a_2},
\qquad
\cos \theta = \frac{h}{R} = -\frac{a_1}{2\sqrt{a_2}},
\qquad
\sin \theta = \frac{v}{R} = \sqrt{1-\frac{a_1^2}{4a_2}}.
\]

Hence
\[
y^c = R^t\left(A_5\cos \theta t + A_6\sin \theta t\right),
\]
where
\[
R=\sqrt{a_2},
\qquad
\theta \text{ is defined as }
\begin{cases}
\cos\theta = -\dfrac{a_1}{2\sqrt{a_2}},\\[0.8em]
\sin\theta = \sqrt{1-\dfrac{a_1^2}{4a_2}}.
\end{cases}
\]

\section{Convergence under higher-order difference equation}

\begin{enumerate}[label=\arabic*$^\circ$]
    \item Distinct root: $b_{\mathrm{dom}}$ is the root with highest absolute value.
    Convergence requires
    \[
    \left|b_{\mathrm{dom}}\right|<1.
    \]

    \item Repeated root:
    \[
    \left|b\right|<1 \Rightarrow \text{converges}.
    \]

    \item Complex root:
    \[
    R>1: \text{ explosive oscillations},
    \qquad
    R<1: \text{ damped oscillations}.
    \]
\end{enumerate}

\section{Simultaneous difference equation (state-space form)}

Let
\[
x_t=y_{t+1}.
\]

Then
\[
\begin{cases}
x_{t+1}+a_1x_t+a_2y_t=C,\\
x_t=y_{t+1}.
\end{cases}
\]

Therefore,
\[
\begin{pmatrix}
x_{t+1}\\
y_{t+1}
\end{pmatrix}
=
\begin{pmatrix}
-a_1 & -a_2\\
1 & 0
\end{pmatrix}
\begin{pmatrix}
x_t\\
y_t
\end{pmatrix}
+
\begin{pmatrix}
C\\
0
\end{pmatrix}.
\]

Generalize:
\[
X_{t+1}=AX_t+B
\]
where
\[
X_t=\left(x_{1t},x_{2t},\dots,x_{nt}\right)
\]
is state space.

\[
A \text{ is a constant matrix with } a_{ij} \text{ as elements } \left(i,j=1,\dots,n\right),
\]
\[
B \text{ is an array with } B_i,\ i=1,\dots,n \text{ as elements.}
\]

In this example,
\[
X_{t+1}=
\begin{pmatrix}
x_{t+1}\\
y_{t+1}
\end{pmatrix},
\qquad
A=
\begin{pmatrix}
-a_1 & -a_2\\
1 & 0
\end{pmatrix},
\qquad
b=
\begin{pmatrix}
C\\
0
\end{pmatrix}.
\]

Everything about dynamics and stability is controlled by eigenvalues of $A$.

Ignore constant term $B$ to study deviation from steady state:
\[
X_{t+1}=AX_t.
\]

Try solution:
\[
X_t = vb^t.
\]

Plug in:
\[
vb^{t+1}=A\left(vb^t\right)
\quad\Rightarrow\quad
bv=Av.
\]

Here $b$ is an eigenvalue of $A$, with eigenvector $v$.

A nonzero $v$ exists iff
\[
\det\left(A-bI\right)=0.
\]

\section{Characteristic equation}

Characteristic equation:
\[
p\left(b\right)=\left|A-bI\right|=0.
\]

For a $2\times 2$ matrix:
\[
p\left(b\right)=b^2-Tb+D=0,
\]
where
\[
T=\mathrm{trace}\left(A\right),\qquad D=\det\left(A\right).
\]

Also,
\[
T=b_1+b_2,\qquad D=b_1b_2,
\]
so
\[
p\left(b\right)=b^2-\left(b_1+b_2\right)b+b_1b_2=\left(b-b_1\right)\left(b-b_2\right).
\]

Eigenvalues are the solution of the characteristic equation: $b_1,b_2$.

Back to the example, since
\[
A=
\begin{pmatrix}
-a_1 & -a_2\\
1 & 0
\end{pmatrix},
\]
we have
\[
T=\mathrm{tr}\left(A\right)=-a_1,
\qquad
D=\det\left(A\right)=a_2.
\]

Thus characteristic equation:
\[
b^2+a_1b+a_2=0.
\]

Eigenvalues:
\[
\frac{-a_1\pm \sqrt{a_1^2-4a_2}}{2},
\qquad \text{if } a_1^2>4a_2.
\]

\section{\texorpdfstring{Stability triangle on the $\left(T,D\right)$ plane}{Stability triangle on the (T,D) plane}}

\begin{enumerate}[label=\arabic*$^\circ$]
    \item
    \[
    \Delta=T^2-4D.
    \]
    \[
    \begin{cases}
    \Delta<0: & \text{complex roots } \Rightarrow \text{above red curve},\\
    \Delta>0: & \text{real roots } \Rightarrow \text{below red curve}.
    \end{cases}
    \]

    \item
    \[
    p\left(1\right)=0,\quad p\left(-1\right)=0
    \]
    are the black curves.

    \[
    p\left(1\right)=1-T+D,
    \qquad
    \text{right to } p\left(1\right)=0 \Rightarrow p\left(1\right)<0.
    \]

    \[
    p\left(-1\right)=1+T+D,
    \qquad
    \text{left to } p\left(-1\right)=0 \Rightarrow p\left(-1\right)<0.
    \]

    \item
    \[
    D=1 \text{ and } \left|b\right|=1 \text{ for complex roots}
    \]
    is the green curve.
\end{enumerate}
\begin{center}
    \IfFileExists{\ProjectRoot/_assets/figures/T-D plane.png}{%
        \includegraphics[width=0.5\linewidth]{T-D plane.png}%
    }{%
        \fbox{\texttt{missing figure: T-D plane.png}}%
    }
\end{center}
Why $D$ controls $\left|b\right|$ in complex case?

If roots are complex conjugates
\[
b_{1,2}=re^{\pm i\theta}
\quad\Rightarrow\quad
D=b_1b_2=r^2
\quad\Rightarrow\quad
r=\sqrt D.
\]

Stable ($r<1$) implies
\[
D<1.
\]

Thus, above green curve $\left|b\right|>1$, below green curve $\left|b\right|<1$.

\textcolor{red}{
* Conditions for two stable roots (both inside unit circle)
}

\[
D<1
\quad\text{and}\quad
p\left(1\right)>0
\quad\text{and}\quad
p\left(-1\right)>0
\]
\[
\Rightarrow\quad
D<1,\quad
1-T+D>0,\quad
1+T+D>0
\]
\[
\Rightarrow\quad
D<1
\quad\text{and}\quad
\left|T\right|<1+D.
\]

\[
\begin{cases}
\text{complex roots with modulus }<1, & \text{if }\Delta<0 \ \left(T^2<4D\right),\\
\text{real roots both between }-1\text{ and }1, & \text{if }\Delta>0 \ \left(T^2>4D\right).
\end{cases}
\]

\textcolor{red}{
* Conditions for one stable root and one unstable root.
}

\[
p\left(1\right)<0 \ \&\ p\left(-1\right)>0
\qquad \text{or} \qquad
p\left(1\right)>0 \ \&\ p\left(-1\right)<0.
\]

That is,
\[
\begin{cases}
1-T+D<0,\\
1+T+D>0,
\end{cases}
\qquad \text{or} \qquad
\begin{cases}
1-T+D>0,\\
1+T+D<0.
\end{cases}
\]

Thus,
\[
\left|T\right|>\left|1+D\right|.
\]

\[
p\left(1\right)\ \&\ p\left(-1\right)
\text{ has opposite signs }
\Rightarrow
\text{There is a root in the interval }\left(-1,1\right).
\]

Since
\[
p\left(1\right)<0,\ p\left(b\right)\to +\infty \text{ as } b\to +\infty
\]
or
\[
p\left(-1\right)<0,\ p\left(b\right)\to +\infty \text{ as } b\to -\infty,
\]
there is another root outside the interval $\left[-1,1\right]$.

\textcolor{red}{
* Conditions for explosive and unstable roots (both roots outside unit circle or complex with modulus $>1$)
}

Four cases:
\begin{enumerate}[label=\arabic*$^\circ$]
    \item $p\left(1\right)<0,\ p\left(-1\right)<0,\ D<0,\ \left|T\right|>1$
    \item $p\left(1\right)>0,\ p\left(-1\right)>0,\ D>1,\ T<-2$
    \item $\Delta<0$ and $D>1$
    \item $p\left(1\right)>0,\ p\left(-1\right)>0,\ D>1,\ T>2$
\end{enumerate}

\section{Region-by-region summary}

\[
\begin{array}{|c|c|c|}
\hline
\text{Region} & \text{Conditions} & \text{Roots} \\
\hline
1 & p\left(1\right)<0 \ \&\ p\left(-1\right)>0 &
\text{saddle: }
\begin{cases}
1 \text{ stable root in }\left(-1,1\right),\\
1 \text{ unstable root in }\left(1,+\infty\right)
\end{cases}
\\
\hline
2 & p\left(1\right)<0 \ \&\ p\left(-1\right)<0 &
\text{explosive: }
\begin{cases}
1 \text{ unstable root in }\left(1,+\infty\right),\\
1 \text{ unstable root in }\left(-\infty,-1\right)
\end{cases}
\\
\hline
3 & p\left(1\right)>0 \ \&\ p\left(-1\right)<0 &
\text{saddle: }
\begin{cases}
1 \text{ stable root in }\left(-1,1\right),\\
1 \text{ unstable root in }\left(-\infty,-1\right)
\end{cases}
\\
\hline
4 &
\begin{array}{c}
p\left(1\right)>0 \ \&\ p\left(-1\right)>0 \ \& \\
\Delta>0 \ \&\ D>1 \ \left(T<-2\right)
\end{array}
&
\text{explosive: 2 real unstable roots in }\left(-\infty,-1\right)
\\
\hline
5 &
\begin{array}{c}
p\left(1\right)>0 \ \&\ p\left(-1\right)>0 \ \& \\
\Delta<0 \ \&\ D>1
\end{array}
&
\begin{array}{c}
\text{explosive: complex conjugates with} \\
\text{modulus } r=\sqrt D>1,\ \text{oscillates \& diverges}
\end{array}
\\
\hline
6 &
\begin{array}{c}
p\left(1\right)>0 \ \&\ p\left(-1\right)>0 \ \& \\
\Delta<0 \ \&\ D<1 \ \left(-2<T<2\right)
\end{array}
&
\begin{array}{c}
\text{stable: complex conjugates with} \\
\text{modulus } r=\sqrt D<1,\ \text{oscillates \& converges}
\end{array}
\\
\hline
7 &
\begin{array}{c}
\Delta>0 \ \&\ D<1 \ \&\ p\left(1\right)>0 \ \&\ p\left(-1\right)>0
\end{array}
&
\begin{array}{c}
\text{stable: 2 real roots } \left|b_1\right|<1,\ \left|b_2\right|<1
\end{array}
\\
\hline
7a &
D>0 \ \&\ T<0 &
\begin{array}{c}
-1<b_2\le b_1<0 \\
\text{\small (same sign, and }T=b_1+b_2\text{ negative)} \\
\text{\small decay but can flip sign}
\end{array}
\\
\hline
7b &
D<0 &
\begin{array}{c}
-1<b_2<0<b_1<1 \\
\text{\small (opposite sign, }D=b_1b_2)
\text{\small decay but can flip sign}
\end{array}
\\
\hline
7c &
D>0 \ \&\ T>0 &
\begin{array}{c}
0<b_2\le b_1<1 \\
\text{\small (same sign, and }T=b_1+b_2\text{ positive)} \\
\text{\small monotone converge}
\end{array}
\\
\hline
8 &
\begin{array}{c}
\Delta>0,\ D>1,\ T>2 \\
p\left(1\right)>0,\ p\left(-1\right)>0
\end{array}
&
\begin{array}{c}
\text{explosive: 2 real roots} \\
b_1>1,\ b_2>1
\end{array}
\\
\hline
\end{array}
\]

\section{Blanchard and Kahn (BK) conditions}

Predetermined variable: initial value is given, states like $k,h,b$ {\small\textcolor{black}{(i.e.\ capital, hour, bond)}}.

Jump variable: initial value not given, can jump to satisfy equilibrium, like $c,p$.

\[
\text{Stable roots} = \#\text{ predetermined variables}
\Rightarrow
\text{Saddle path (Determinacy)}
\]

\[
\text{Stable roots} < \#\text{ predetermined variables}
\Rightarrow
\text{Source (Explosive, no convergent equilibrium)}
\]

\[
\text{Stable roots} > \#\text{ predetermined variables}
\Rightarrow
\text{Sink (Indeterminacy)}
\]

\section{Why?}

Linearize around steady state, we write deviations as a first-order system:
\[
X_{t+1}=AX_t.
\]

If $A$ has eigenpairs $\left(b_i,v_i\right)$, the general solution can be written as sum of modes:
\[
X_t=\sum_i \alpha_i b_i^t v_i.
\]

If
\[
\left|b_i\right|<1,
\]
mode dies out $\Rightarrow$ stable.

If
\[
\left|b_i\right|>1,
\]
mode blows up $\Rightarrow$ unstable / explosive.

\[
\text{A convergent equilibrium path must impose: }
\alpha_i=0
\quad\text{for every } \left|b_i\right|>1,
\]
which is a set of restrictions, one restriction per unstable root.

Now partition variables into
\[
\left\{\text{predetermined states } s_t \ \left(\text{given } s_0\right)\ \text{with dimension } n_s,\right.
\]
\[
\left.\text{jump variables } j_t \ \left(\text{free } j_0\right)\ \text{with dimension } n_j\right\}.
\]

\[
X_t=\left(s_t,j_t\right).
\]

Initial conditions pin down $s_0$, but can choose $j_0$.

Let
\[
n_u=\#\text{ unstable roots},
\qquad
n_\ell=\#\text{ stable roots}.
\]

\[
n_u+n_\ell=n_s+n_j.
\]

To satisfy $n_u$ restrictions, we have $n_j$ free choices to satisfy them. Thus
\[
n_u=n_j
\qquad\Longleftrightarrow\qquad
n_\ell=n_s.
\]

\subsection{Determinacy}
\[
n_\ell=n_s \quad \left(n_u=n_j\right):
\]
jump variables jump to land on saddle path.

\subsection{No convergent equilibrium}
\[
n_\ell<n_s \quad \left(n_u>n_j\right):
\]
too many unstable modes to eliminate, but too few jump variables to adjust.

\subsection{Indeterminacy}
\[
n_\ell>n_s \quad \left(n_u<n_j\right):
\]
more jump freedom needed to kill unstable modes (multiple equilibria)

\ThisNoteEnd
